\subsection{8. Während der Arbeit entdeckte Änderung: „Kommutant 7``
richtig
einordnen}\label{wuxe4hrend-der-arbeit-entdeckte-uxe4nderung-kommutant-7-richtig-einordnen}

Die strikte Originalquellenprüfung stoppte korrekt, als eine andere
Arbeit den früheren nativen Bericht ergänzte. Der alte Stand mit Hash
\codeword{1062e13fe0fb79cce015657a7b4d40a61e9639922ff95fa60c689c0068c38ffc}
blieb eingefroren. Der gesonderte geänderte Entwurf hat Hash
\codeword{ab745e74b1232b1ddf35dd0e0d0da4ace0176c797d769e5db8bedd0010bcba6a}.
Die in dieser Revision verwendeten Energie- und Polzertifikate änderten
sich bei dieser Prüfung nicht. Der Fehlerbeleg wurde aufbewahrt.

Der neue Entwurf listet drei dunkle Darstellungstypen mit Dimensionen
24.000, 11.200, 2.688; drei helle mit 2.880, 576, 320; und einen χ-Typ
mit 64. \textbf{Unter Voraussetzung dieser Zerlegung} lässt sich seine
Folgerung exakt nachprüfen:

\[\adjustbox{max width=\linewidth}{$\displaystyle 
 \mathcal H_{N=3}\simeq
 \bigoplus_{d\in D}V_d\ \oplus\
 \bigoplus_{b\in B}(\mathbb C^2\otimes V_b)\ \oplus V_\chi.
$}\]

Die Summe ist \(37888+2\cdot3776+64=45504\). Die drei hellen Typen
kommen im \textbf{vollen} Raum jeweils zweimal vor. Daraus folgen:

\begin{itemize}
\tightlist
\item
  Nur die Gruppenwirkung hat einen Kommutanten der Dimension
  \(3+3\cdot4+1=16\), nicht sieben.
\item
  Werden zusätzlich die nichtverschwindenden hellen Paarmischungen und
  \(N_b\) als Operationen zugelassen, erzeugen sie die vollständigen
  \(M_2\)-Multiplizitätsalgebren. Dann bleiben sieben zentrale Skalare.
\item
  Sieben ist \textbf{keine absolute Untergrenze für beliebige weitere
  Operationen}. Ein erlaubter Operator, der diese sieben Blöcke
  verbunden mischt, zwingt die sieben Skalarwerte zur Gleichheit und
  lässt nur einen übrig. Der endliche Inzidenzmatrix-Test hat Rang sechs
  und beweist dieses Gegenargument.
\item
  Die Zerlegung allein beweist weder operative Verfügbarkeit der
  kontinuierlichen Symmetrie noch einen Einzelfermion-Link. Selbst alle
  bankinternen Symmetrieoperationen erhalten die lokale Fermionparität.
\end{itemize}

Der interessante Wert sieben wird dadurch nicht verworfen. Korrigiert
werden die Aussagen „voller N=3-Raum multiplizitätsfrei``, „absolute
Grenze für jede Operation`` und eine unbewiesene Gleichsetzung von
Symmetrie und verfügbarem Instrument. Die genaue Darstellungszerlegung
selbst bleibt in diesem Zusatz eine externe Voraussetzung, kein frisch
reproduzierter Satz.

Der Debugging-Skill führte hier zur Versionsprüfung statt zum Ersetzen
eines Pins. Das neue Manifest trennt deshalb einen erfolgreichen
\textbf{eingefrorenen Quellenreplay} von
\codeword{originals\_unchanged=false}.

\subsection{9. T1-T8: was dieses Ergebnis beiträgt und nicht
beiträgt}\label{t1-t8-was-dieses-ergebnis-beitruxe4gt-und-nicht-beitruxe4gt}

Die Begriffe folgen der aktuellen offenen Problemliste und ihrer
Forschungszuordnung. Keine Akzeptanzmarkierung wurde heraufgesetzt.

{\def\LTcaptype{none} % do not increment counter
\begin{longtable}[]{@{}
  >{\raggedright\arraybackslash}p{(\linewidth - 4\tabcolsep) * \real{0.3333}}
  >{\raggedright\arraybackslash}p{(\linewidth - 4\tabcolsep) * \real{0.3333}}
  >{\raggedright\arraybackslash}p{(\linewidth - 4\tabcolsep) * \real{0.3333}}@{}}
\toprule\noalign{}
\begin{minipage}[b]{\linewidth}\raggedright
Tor
\end{minipage} & \begin{minipage}[b]{\linewidth}\raggedright
Fehlender Nachweis
\end{minipage} & \begin{minipage}[b]{\linewidth}\raggedright
Beitrag dieser Revision
\end{minipage} \\
\midrule\noalign{}
\endhead
\bottomrule\noalign{}
\endlastfoot
T1 & Ursprung von P1/P2, Dimension, Compiler- und Zustandswahl &
Bekannte gerade Operationsalgebra ist für Einzeltransport unzureichend;
ursprüngliche Auswahl weiter offen \\
T2 & Tatsächlich halbgeladenes, markiertes E8-Feld samt Renormierung,
Energie- und Adjungiertenkontrolle & Der Ladung-eins-Lochzustand ist
kontrolliert, aber nicht mit dem fehlenden Half-Charge-Feld
identifiziert \\
T3 & Ein gemeinsamer, aus TFPT ausgewählter lokaler unitärer 3+1D-Parent
& Vollständiger bedingter Zwei-Banken-Transfertest; kein ausgewählter
räumlicher Parent \\
T4 & Chirales SM-Maß, Anomalien/Index, gleichmäßige Spiegelentkopplung &
Feldtyp-Widerspruch weiter sichtbar; kein chirales Maß \\
T5 & Wechselwirkender Kontinuumsübergang, Lorentzverhalten, Clustering,
Confinement und Streuung & Endliche-Sektor-Fehlerkontrolle ist ein
Baustein, aber kein räumlicher Grenzübergang \\
T6 & Interne Herleitung aller Eichkopplungen und vollständiger
Neutrinotextur/-skala & Keine neue Herleitung; \(\tau,\kappa\) sind
deklarierte Testparameter \\
T7 & Masseloser dynamischer Spin zwei, beide Helizitäten, universelle
Kopplung im selben Parent & Offen; der neue pauschale Ausschluss eines
Dimension-vier-(1,1)-Tensors wurde korrigiert \\
T8 & Physische Präparation/Anfangszustand und ein gemeinsames
Quellfunktional aller Auslesungen & Anfänglicher Polfilter wird für
\textgreater89,7 \% unnötig; geladene Präparation und Messinstrument
weiter offen \\
\end{longtable}
}

\subsection{10. Nächste entscheidende Arbeiten mit klaren
Abbruchkriterien}\label{nuxe4chste-entscheidende-arbeiten-mit-klaren-abbruchkriterien}

\textbf{A. Die Quelle des geraden Gesamtlinks finden.} Nicht noch eine
riesige Kontrollmatrix bauen, sondern für jeden tatsächlich
ursprünglichen Generator \(O\) prüfen, ob \([O,\Pi_x]\ne0\) bei
erhaltener Gesamtparität möglich ist. Eine positive Antwort muss
Generator, Teilbereichsdefinition, Hilbertraum, Gaussreferenz und
Matrixelement liefern. Bleiben alle Generatoren lokal gerade, ist die
Suche nach dem Link durch deren weitere Komposition beendet. Dann muss
die ursprüngliche Teilraumdefinition oder der Quellenvertrag explizit
geändert werden; man darf keinen Hopterm als Ergebnis ausgeben.

\textbf{B. Die ursprüngliche Entnahme als gemeinsames Instrument
realisieren.} Eine Referenzmode könnte Ladung erhalten, während
\(f_r\Omega\) präpariert und ausgelesen wird. Zu prüfen sind ihre
Herkunft, Energiekosten, Erfolgshäufigkeit und Erhaltung der internen
Kohärenz. Der neue Satz entfernt bereits die zusätzliche Pflicht eines
perfekten anfänglichen Energiefilters. Er entfernt nicht die Pflicht,
ein geladenes Instrument auszuweisen.

\textbf{C. Den Feldtyp vor räumlichen Rechnungen entscheiden.} Entweder
folgt ein passender Lorentzträger wirklich aus derselben Quelle, oder
der skalare Weyl-Ansatz wird in seiner jetzigen Form verworfen. Eine
bloße zusätzliche Verdopplung ohne Herkunft ist ein Modellvorschlag,
keine Ableitung.

\textbf{D. Erst dann viele Banken.} Der nächste räumliche Test benötigt
dieselbe Quelle, einen bestimmten Graphen, Lokalisierung der Operationen
und mit dem Volumen verträgliche Fehler- und Energiebounds. Ein einziges
neues Linklemma liefert weder Raumdimension noch Lichtkegel, chirale
Materie oder Gravitation.

\textbf{E. Arithmetik getrennt scharf halten.} Eine echte Verbindung
müsste eine gemeinsame Quellabbildung mit Generator,
Zustands-/Spurstruktur, Phasen und Domänen liefern. Die volle signierte
Weilform oder die ganze Determinanten- identität bleibt das Beweisziel.
Der bloße Auftritt von Primzahlen erfüllt keines davon. Der globale
RH-Index ließ sich wegen einer fehlenden historischen Quelle und
veränderter Review-Quellen nicht vollständig aktualisieren; die gezielte
Originalquellenprüfung ist enger als ein frischer Gesamtindex. Es wurde
kein RH-Abschlusskandidat registriert.

\subsection{11. Reproduktion und ehrliche
Liefergrenze}\label{reproduktion-und-ehrliche-liefergrenze}

\codeword{replay.py --frozen-only} führt die fünf Prüfer normal und mit
\codeword{-OO} aus, vergleicht beide JSON-Ergebnisse und die zwei
externen Originalberichte, prüft alle eingefrorenen Quellen und
protokolliert Änderungen der Originalorte gesondert. Ein Fehler führt zu
FAIL. Der zuvor aufgetretene Live-Quellenfehler bleibt in
\codeword{source\_drift\_failure\_receipt.json} sichtbar.

Die abgesicherten Zahlen stehen vollständig rational in
\codeword{two\_bank\_transfer\_normal.json}; Dezimalwerte dienen nur der
Lesbarkeit. Die analytische Beweiskette steht in Abschnitt 4. Das Paket
umfasst die unveränderten externen Arbeiten, deren Programme, die native
v1.6.4-Abhängigkeit und alle neuen Berichte. NumPy, SymPy und für das
ältere native Paket dessen separat dokumentierte Voraussetzungen sind zu
unterscheiden.

Die versionierte ausführliche Markdown-Lieferung enthält zusätzlich den
vollständigen eingefrorenen nativen Herleitungsbericht v1.6.4 als
historischen Anhang. Daneben gibt es ein kurzes Update und eine
bildliche Erklärung. \textbf{Die großen TFPT-Haupt-PDFs und die Webseite
werden in dieser Revision nicht als aktualisiert ausgegeben. Es gab
keinen Commit oder Push.}

\begin{center}\rule{0.5\linewidth}{0.5pt}\end{center}

\section{Nachprüfung der beiden zuletzt eingegangenen
Untersuchungen}\label{nachpruxfcfung-der-beiden-zuletzt-eingegangenen-untersuchungen}

Integrierter Nachtrag zu v1.6.5 · 15. September 2026

Quellen: die Anlagen \codeword{8e110ae1-19e2-4d47-bfc8-ebe4fe9ab81b} und
\codeword{b928b72a-6669-4c67-a75d-124095791c6f}, unverändert
eingefroren. Zusätzlich wurden die vier genannten Programme
\codeword{t1\_fixed.py}, \codeword{stabilizer.py},
\codeword{t5\_twobank.py}, \codeword{t5\_varresp.py} vollständig gelesen
und archiviert. Eine eigene Nachrechnung prüft die kritischen
Folgerungen; sie ist \textbf{kein vollständiger Replay aller sechs
Sonden} des fremden Arbeitsordners.

\subsection{1. Übersicht: übernehmen, korrigieren oder
offenlassen}\label{uxfcbersicht-uxfcbernehmen-korrigieren-oder-offenlassen}

{\def\LTcaptype{none} % do not increment counter
\begin{longtable}[]{@{}
  >{\raggedright\arraybackslash}p{(\linewidth - 2\tabcolsep) * \real{0.5000}}
  >{\raggedright\arraybackslash}p{(\linewidth - 2\tabcolsep) * \real{0.5000}}@{}}
\toprule\noalign{}
\begin{minipage}[b]{\linewidth}\raggedright
Neue Aussage
\end{minipage} & \begin{minipage}[b]{\linewidth}\raggedright
Ergebnis der Prüfung
\end{minipage} \\
\midrule\noalign{}
\endhead
\bottomrule\noalign{}
\endlastfoot
Innerer Viererzyklus erhält alle 60 W-Kanäle & Bestätigt, jetzt
einschließlich der richtigen Fermion-Vorzeichen und explizitem
Bosonlift \\
Grundzustandseindeutigkeit/Singulett/Gaplücke weiterhin unbekannt &
Überholter Stand: unter dem festgelegten nativen Vertrag bereits
bewiesen; voller Vektor weiter unbekannt \\
Leerer Grundzustand für alle μ≥0 & Falsch; bei μ=0 liegt der eindeutige
N=64-Grundzustand unter −1,129636Δ \\
Kommutant 7 als absolute Untergrenze & Nur im benannten blocktreuen
erweiterten Operationsvertrag; siehe Hauptbericht Abschnitt 8 \\
Beliebige Einteilchenmatrix ist damit als natives Fock-Wort verfügbar &
Nicht gezeigt; assoziative Matrixalgebra, Lie-Kontrolle und Fock-Lift
sind verschiedene Fragen \\
Kanal-Variationszustand hat E=−0,0196Δ & Die angegebene Phase im Code
ergibt gegen den gepinnten P=f\_jf\_i-Vertrag +0,05736Δ; ein relatives
Vorzeichen korrigiert es \\
99,892-\%-Transfer bei t≈2484 & Numerisches Maximum eines
Vierzustands-Paarmodells im untersuchten Zeitfenster; kein
erster/globaler exakter Maximalsatz und kein nativer Einloch-Transfer \\
Gleiche Zweigableitungen genau bei ε=0 & Algebraisch richtig; die
Eigenwertlücke bleibt dort √32· \\
SU(4)³-Anomalie =16 & Richtig unter A(4)=1 und der zusätzlichen
Interpretation aller (16,4) als gleichhändige Weylfelder; Eichung ist
eine weitere Voraussetzung \\
Gravitationsanomalie =64 & Keine reine perturbative Gravitationsanomalie
in 3+1D; 64 kann bei zusätzlich deklarierter gleichgeladener U(1) die
gemischte U(1)-Gravitationsanomalie zählen \\
Mit einer Ableitung kein (1,1)-Tensor bei Dimension ≤4 & Falsch: der
Energie-Impuls-Tensor ist ein Gegenbeispiel; daraus folgt aber noch kein
dynamisches Graviton \\
Kein lokaler kovarianter kinetischer Term für (1,0) & Zu pauschal: der
bilineare Ein-Ableitungs-Term fehlt, ein Zwei-Ableitungs-Skalar
existiert; Positivität/Constraints/Herkunft bleiben zu prüfen \\
\end{longtable}
}

„Negativ geschlossene Route`` bedeutet nicht „T2, T3, T4 oder T7
geschlossen``. Ein ausgeschlossener Ansatz beantwortet nicht die
jeweilige physikalische Existenzfrage. Die beiden gelieferten Texte
widersprechen sich vor allem beim bereits bewiesenen modellinternen
Grundzustand; sie dürfen nicht als gleichzeitig aktueller einheitlicher
Status zitiert werden.

\subsection{2. Positiver neuer Anschluss: der korrekt angehobene
Viererzyklus}\label{positiver-neuer-anschluss-der-korrekt-angehobene-viererzyklus}

Die Abbildung der Fermionmarken lautet \(p(4s+a)=4s+(a+1\bmod4)\). Auf
Paaren ist zwingend das Exteriorvorzeichen zu beachten: Falls
\(p(i)>p(j)\), erhält das sortierte Paar ein Minuszeichen. Der
gelieferte \codeword{stabilizer.py} lässt dieses Zeichen weg. Das wäre
im Allgemeinen falsch; für den untersuchten Viererzyklus überlebt das
positive Resultat jedoch auch den korrekten Test.

Alle 60 W-Zeilen werden mit ihren Vorzeichen wieder auf W-Zeilen
abgebildet. Wir haben den zugehörigen signierten Bosonoperator \(R_b\)
explizit gebildet und exakt geprüft:

\[\adjustbox{max width=\linewidth}{$\displaystyle 
 W'=R_bW,\qquad R_bR_b^T=I_{60},\qquad R_b^4=I_{60},
 \qquad R_b^2\ne I_{60}.
$}\]

Damit existiert ein \textbf{gemeinsamer nativer Tensor-Automorphismus}
auf Fermionen und Bosonen, nicht nur eine Übereinstimmung von
Zeilenzahlen. Die gleichzeitige Transformation erhält Paarwechselwirkung
und Bosonzahl. Sie beweist eine Ordnung-vier-Symmetrie, nicht deren
operative Verfügbarkeit.

Dieser Permutationszyklus ist nicht die zentrale Matrix \(iI_4\) von
SU(4). Seine Determinante auf dem Viererfaktor ist −1. Mit einer
zusätzlichen Ladungsphase kann man geeignete Gruppenlifts vergleichen;
die Permutation, die zentrale Phase, der Clock und die geometrische
Glue-Markierung sind dadurch aber nicht schon identifiziert. Genau diese
Intertwiner-Frage ist ein sinnvoller neuer Anschluss, statt noch einmal
nur vier zu zählen.

\subsection{3. Warum der Einteilchen-Abschluss den Fock-Operationssatz
nicht
schließt}\label{warum-der-einteilchen-abschluss-den-fock-operationssatz-nicht-schlieuxdft}

Die Irreduzibilität der Einteilchendarstellung kann ihre assoziative
Matrixalgebra zu \(B(\mathbb C^{64})\) machen. Daraus folgt nicht, dass
jede Matrix als ein zulässiges physisches Wort verfügbar ist, und auch
nicht, dass ihre zweite Quantisierung bereits erzeugt wird.

Ein kleinstes exaktes Gegenbeispiel zur falschen Liftregel:
\(A=|1\rangle\langle1|\), \(B=|2\rangle\langle2|\) auf zwei Moden. Dann
\(AB=0\), also \(d\Gamma(AB)=0\), aber
\(d\Gamma(A)d\Gamma(B)=n_1n_2\ne0\). Der zweite Quantisierungsschritt
ist ein Lie-, nicht ein assoziativer Algebra-Homomorphismus dieser Art.

Auch \(\operatorname{diag}(i,1,1,1)\) und
\(\operatorname{diag}(i,i,1,1)\) haben Determinanten i beziehungsweise
−1. Ihre Zugehörigkeit zur \textbf{linearen Matrixalgebra} aus
SU(4)-Generatoren und Identität ist kein exaktes SU(4)-Gruppenwort.
Projektive Gleichheit oder eine zusätzliche U(1)-Phase kann helfen, muss
aber mit der Bosonwirkung, Ladung und Referenz gemeinsam ausgewiesen
werden.

Der fremde Code folgert außerdem die Irreduzibilität des
37.888-dimensionalen dunklen Dreifermionraums aus der vollen
Einteilchenmatrixalgebra. Diese Schlussregel ist nicht begründet.
Bereits die im anderen neuen Text angegebenen drei dunklen irreduziblen
Typen widersprechen der pauschalen Eins-Block-Lesart unter der bloßen
Quellsymmetrie. Die Zahl 8.732.673 mag als sehr schwache obere Schranke
anderweitig verträglich sein; die im Code angegebene Herleitung belegt
sie nicht. Die SVD-/Toleranzprüfungen des Codes sind zudem numerisch,
nicht allein wegen „PASS`` exakte Lie-Beweise.

\subsection{4. Zustandswahl: keine Rückkehr vor den abgesicherten
Grundsatz}\label{zustandswahl-keine-ruxfcckkehr-vor-den-abgesicherten-grundsatz}

Der vorhandene native Satz beweist Eindeutigkeit, N=64, Singulett und
positive Lücke bei μ=0 im angegebenen Kopplungsbereich. Eine komplette
Clebsch-Gordan-Serie oder Voll-Diagonalisierung ist dafür nicht nötig;
Vergleichsungleichungen und Minmax waren gerade der einfache Ausweg. Die
physische Auswahl von H beziehungsweise μ=0 bleibt eine andere Frage.

Das grobe Sandwich \([-1.2,-.0196]\Delta\) ist nach einer
Phasenkorrektur zwar verträglich, aber wesentlich schwächer als
\((-1.158089,-1.129636)\Delta\). Es ist keine Verschärfung.

Im tatsächlich angegebenen Variationscode werden die Fermionen in der
Reihenfolge j, dann i gelöscht. Das liefert \(f_if_jF=-f_jf_iF\),
während der native Vertrag \(P_A=\sum W_{A,ij}f_jf_i\) verwendet. Bei
unverändertem \(g=+\Delta/20\) und den angegebenen
Variationskoeffizienten folgt deshalb

\[\adjustbox{max width=\linewidth}{$\displaystyle 
 E_{\rm Code}/\Delta=\frac12-\frac{23\sqrt3}{90}
 =+.057364793621\ldots,
$}\]

nicht der behauptete negative Wert. Mit korrigierter relativer Phase
lautet er \(1/2-3\sqrt3/10=-.019615242271\ldots\). Alternativ wäre eine
konsequente andere P/g-Phasenkonvention möglich; die Quelle muss sie
dann überall führen. Die gleichzeitige Einteilchendichte des einfachen
Variationszustands bleibt von dieser Phasenreparatur unberührt, weil
seine Bosonzahlkomponenten orthogonal sind. Sie ist keine dynamische
Greenfunktion auf dem nativen Grundzustand und ersetzt dessen schon
kontrollierte Antwort nicht.

\subsection{5. Feldtheorie: falsche Verbote entfernen, echte Bedingungen
behalten}\label{feldtheorie-falsche-verbote-entfernen-echte-bedingungen-behalten}

